\subsection{Query Rewriting}
\label{sec:query_rewriting}

The key way ProvSQL implements the semantics of the extended algebra over
annotated relations is by \emph{rewriting the query} over annotated
relations to include in the query the necessary operations on provenance
operations. The query can then be evaluated using the standard query
evaluator of PostgreSQL. We
% lax begin Lax392996.RewritingRules
employ the following \ref{rn} rewriting rules, which are applied inductively on an
extended relational algebra formula:
\lean{QueryRewriting}{query.Rewriting}
\begin{asparaenum}[(R1)]
	\item \label{r1} \label{rprojection} for $k\in\bN$, $q\in\RA_{k}$,
	$t_1,\dots, t_n$ terms of max-index~\mbox{$\leq k$},
	$\Pi_{t_1,\dots,t_n}(q)$ is rewitten to $\Pi_{t_1,\dots,t_n,\#(k+1)}(q)$
	\item \label{rcp}
	For $k_1,k_2\in\mathbb{N}$, $q_1\in\RA_{k_1}$, $q_2\in\RA_{k_2}$,
	\(q_1\times q_2\) is rewritten to:
	\(\Pi_{\#1,\dots,\#k_1,\#(k_1+2),\dots,\#(k_1+k_2+1),\#(k_1+1)\otimes\#(k_1+k_2+2)}(q_1\times
	q_2)\).
	\item \label{rdupelim} For $k\in\mathbb{N}$, $q\in\RA_k$,
	$\epsilon(q)$ is rewritten to:
	$\gamma_{1,\dots,k}[\#(k+1): \bigoplus](q)$.
	\item For $k\in\mathbb{N}, q_1,q_2\in\RA_k$, $q_1-q_2$ is rewritten to:\\
    \hspace*{1.2em}$\Pi_{\#1,\dots,\#(k+1)}(q_1\bowtie_{\scriptscriptstyle\#1=\#(k+2)\land\dots\land\#k=\#(2k+1)}$\\\hspace*{4em} $\epsilon(\Pi_{\scriptscriptstyle\#1,\dots,\#k}(q_1)-\Pi_{\scriptscriptstyle\#1,\dots,\#k}(q_2)))$\\
	\({}\uplus
  \Pi_{\#1,\dots,\#k,\#(k+1)\ominus\#(2k+2)}\!(q_1\bowtie_{\scriptscriptstyle\#1=\#(k+2)\land\dots\land\#k=\#(2k+1)}\)\\\hspace*{4em}\(
	\gamma_{\scriptscriptstyle\#1,\dots,\#k}[{\scriptstyle\#(k+1):\bigoplus}](q_2))
	\)
% lax end
	\item For $k\in\mathbb{N}$, $q\in\RA_k$, distinct $(i_j)_{1\leq j\leq k}$, terms
	$t_1,\dots,t_n$ of max-index $\leq k$, monoid aggregate functions
	$f_1,\dots,f_n$, $\gamma_{i_1,\dots,i_m}[t_1:f_1,\dots,t_n:f_n](q)$ is
	rewritten to:\\
	$\gamma_{i_1,\dots,i_m}[t_1*\#(k+1):\hat f_1,\dots,
		t_n*\#(k+1):\hat f_n,\#(k+1):\delta\left(\bigoplus\right)](q)\).
	\label{rn}
\end{asparaenum}

Note that
relation names, selection, and multiset sum operators
are not rewritten. \revision{Using query rewriting to perform provenance evaluation
has been proposed in previous works (see, in particular,
Amsterdamer et al. \cite{amsterdamer2011provenance}, Fink and Olteanu
\cite{fink2012aggregation} and Perm \cite{DBLP:conf/icde/GlavicA09}).
The two distinguishing aspects of this work is that
we natively support the multiset semantics that SQL engines natively
implement (requiring adaptation of rewriting rules for projection and
explicit duplicate elimination) and that we support difference.}

We show that these rewriting rules allow us to recover the exact
semantics of the extended relational algebra over annotated relations:

% lax begin Lax392996.RewritingCorrectness
\begin{theoremrep}
  \label{thm:rewriting} \lean{QueryRewriting}{Query.rewriting_valid}
	Let $\mathcal{D}$ be a database schema, $q$ any extended relational
	algebra query over~$\mathcal{D}$, $\mathbb{K}$ an appropriate algebraic
	structure, and $\hat I$ a $\mathbb{K}$-instance over $\mathcal{D}$. Let
	$\hat q$ be the query rewritten from $q$ by applying the rewriting rules
	(R\ref{r1})--(R\ref{rn}) recursively bottom up. Then $\ansem{\hat
			I}q=\sem{\hat I}{\hat q}$.
\end{theoremrep}
% lax end

\begin{proof}
% lax begin Lax392996Proofs.Bridge.rewriting_valid
	We proceed by induction on the structure of~$q$.

  \begin{itemize}
		\item \textbf{relation} for any relation label~$R$ in the domain
			of~$\mathcal{D}$,
			\(
			\sem{\hat I}{\hat R}=\sem{\hat I}{R}=\hat I(R)=\ansem{\hat I}{R}
			\).

		\item \textbf{projection} for $k\in\bN$, $q\in\RA_{k}$,
			$t_1,\dots, t_n$ of terms of max-index $\leq k$,
      \begin{align*}
				\sem{\hat I}{\widehat{\Pi_{t_1,\dots,t_n}(q)}}
				=\sem{\hat I}{\Pi_{t_1,\dots,t_n,\#(k+1)}(\hat q)}
        =&\mset{(t_1(u),\dots,t_n(u),\alpha)\mid (u,\alpha)\in \sem{\hat
        I}{\hat q}}\\
        &=\mset{(t_1(u),\dots,t_n(u),\alpha)\mid (u,\alpha)\in \ansem{\hat I}{q}}
				=\ansem{\hat I}{\Pi_{t_1,\dots,t_n}(q)}.
      \end{align*}
		\item \textbf{selection} for $k\in\bN$, $q\in\RA_k$, and
			$\phi$ a Boolean combination of (in)equality comparisons involving terms
			of max-index $\leq k$,
			\[
				\sem{\hat I}{\widehat{\sigma_\phi(q)}}
				=\sem{\hat I}{\sigma(\phi)(\hat q)}
				=\mset{u'\mid u'\in \sem{\hat I}{\hat q}, \phi(u')}
				=\mset{(u,\alpha)\mid (u,\alpha)\in \ansem{\hat I}{q}, \phi(u)}
				=\ansem{\hat I}{\sigma_\phi(q)}.
			\]
			Indeed, since $\phi$ does not involve ``$\#(k+1)$'',
			$\phi(u')=\phi((u,\alpha))=\phi(u)$.

		\item \textbf{cross product} for $k_1,k_2\in\bN$, $q_1\in\RA_{k_1}$,
			$q_2\in\RA_{k_2}$,
			\begin{align*}
				\sem{\hat I}{\widehat{q_1\times q_2}}
				 & =\sem{\hat I}{\Pi_{\#1,\dots,\#k_1,\#(k_1+2),\dots,\#(k_1+k_2+1),\#(k_1+1)\otimes\#(k_1+k_2+2)}(\hat
				q_1\times\hat q_2)}                                                                                     \\
				 & =\mset{(u_1,\dots,u_{k_1},v_1,\dots,v_{k_2},\alpha\otimes\beta)\mid
				(u_1,\dots,u_{k_1},\alpha,v_1,\dots,v_{k_2},\beta)\in \sem{\hat I}{\hat q_1\times\hat q_2}}             \\
				 & =\mset{(u_1,\dots,u_{k_1},v_1,\dots,v_{k_2},\alpha\otimes\beta)\mid
				(u_1,\dots,u_{k_1},\alpha,v_1,\dots,v_{k_2},\beta)\in \sem{\hat
				I}{\hat q_1}\times\sem{\hat I}{\hat q_2}}                                                               \\
				 & =\mset{(u_1,\dots,u_{k_1},v_1,\dots,v_{k_2},\alpha\otimes\beta)\mid
				(u_1,\dots,u_{k_1},\alpha,v_1,\dots,v_{k_2},\beta)\in \ansem{\hat
				I}{q_1}\times \ansem{\hat I}{q_2}}                                                                      \\
				 & =\ansem{\hat I}{q_1\times q_2}.
			\end{align*}

		\item \textbf{multiset sum} for $k\in\bN$, $q_1,q_2\in\RA_k$,
			\[\sem{\hat I}{\widehat{q_1\uplus q_2}}=\sem{\hat I}{\hat q_1\uplus
					\hat q_2}=\sem{\hat I}{\hat q_1}\uplus\sem{\hat I}{\hat q_2}=
				\ansem{\hat I}{q_1}\uplus\ansem{\hat I}{q_2}=\ansem{\hat I}{q_1\uplus
					q_2}.\]

		\item \textbf{duplicate elimination} for $k\in\bN$,
			$q\in\RA_k$,\\
			\begin{align*}
				\sem{\hat I}{\widehat{\epsilon(q)}}
				 & =
				\sem{\hat I}{\gamma_{1,\dots,k}[\#(k+1):\bigoplus](\hat q)}                   \\
				 & =\set{\left(v_1,\dots,v_k,\bigoplus\mset{\alpha\mid u\in \sem{\hat I}{\hat
							q}, u=(v_1,\dots,v_k,\alpha)}\right)\:\middle|\: (v_1,\dots,v_k)\in \sem{\hat
				I}{\epsilon(\Pi_{1,\dots,k}(\hat q))}}                                        \\
				 & =\set{\left(v_1,\dots,v_k,\bigoplus\mset{\alpha\mid u\in \sem{\hat I}{\hat
							q}, u=(v_1,\dots,v_k,\alpha)}\right)\:\middle|\:
					\sem{\hat I}{\Pi_{1,\dots,k}(\hat
				q))}(v_1,\dots,v_k)>0}                                                        \\
				 & =\set{\left(v_1,\dots,v_k,\bigoplus\mset{\alpha\mid u\in \sem{\hat I}{\hat
							q}, u=(v_1,\dots,v_k,\alpha)}\right)\:\middle|\:
				\exists\alpha\: \sem{\hat I}{\hat q}(v_1,\dots,v_k,\alpha)>0}                 \\
				 & =\set{\left(v_1,\dots,v_k,\bigoplus\mset{\alpha\mid u\in \ansem{\hat I}{
							q}, u=(v_1,\dots,v_k,\alpha)}\right)\:\middle|\:
				\exists\alpha\: \ansem{\hat I}{q}(v_1,\dots,v_k,\alpha)>0}                    \\
				 & = \bigcup_{u\mid\exists \alpha(u,\alpha) (u,va)\in\ansem{\hat I}{q}}
				\set{\left(u,\bigoplus_{\alpha\mid (u,\alpha)\in \ansem{\hat
				I}{q}}\alpha\right)}                                                          \\
				 & =\ansem{\hat I}{\epsilon(q)}.
			\end{align*}

		\item \textbf{multiset difference}
			for $k\in\bN$, $q_1,q_2\in\RA_k$,
			\begin{align*}
        \hspace*{-2em}\sem{\hat I}{\widehat{q_1-q_2}}
				       & =  \lBrack       &            &
				\Pi_{\#1,\dots,\#(k+1)}(\hat
				q_1\bowtie_{\#1=\#(k+1)\land\dots\land\#k=\#(2k)}\epsilon(\Pi_{\#1,\dots,\#k}(\hat
				q_1)-\Pi_{\#1,\dots,\#k}(\hat q_2)))                              \\
				       &                  & {}\uplus{} &
				\Pi_{\#1,\dots,\#k,\#(k+1)\ominus\#(2k+2)}(\hat
				q_1\bowtie_{\#1=\#(k+2)\land\dots\land\#k=\#(2k+1)}
				\gamma_{\#1,\dots,\#k}[\#(k+1):\bigoplus](\hat
				q_2))) & \rBrack^{\hat I}                                         \\
				       & =                &            & \sem{\hat I}{
				\Pi_{\#1,\dots,\#(k+1)}(\hat
				q_1\bowtie_{\#1=\#(k+1)\land\dots\land\#k=\#(2k)}\epsilon(\Pi_{\#1,\dots,\#k}(\hat
				q_1)-\Pi_{\#1,\dots,\#k}(\hat q_2)))}                             \\
				       &                  & {}\uplus{} &
				\sem{\hat I}{\Pi_{\#1,\dots,\#k,\#(k+1)\ominus\#(2k+2)}(\hat
					q_1\bowtie_{\#1=\#(k+2)\land\dots\land\#k=\#(2k+1)}
					\gamma_{\#1,\dots,\#k}[\#(k+1):\bigoplus](\hat
				q_2))}                                                            \\
				       & =                &            &
				\mset{(u,\alpha)\mid (u,\alpha)\in \sem{\hat
						I}{\hat q_1}, u\in \sem{\hat I}{
						\Pi_{\#1,\dots,\#k}(\hat q_1)-
						\Pi_{\#1,\dots,\#k}(\hat q_2)
				}}                                                                \\
				       &                  & {}\uplus{} &
				\mset{(u,\alpha\ominus\beta)\mid (u,\alpha)\in\sem{\hat I}{\hat q_1},
					(u,\beta)\in
					\sem{\hat I}{
						\gamma_{\#1,\dots,\#k}[\#(k+1):\bigoplus](\hat
				q_2))} }                                                          \\
				       & =                &            &
				\mset{(u,\alpha)\mid (u,\alpha)\in \sem{\hat
						I}{\hat q_1}, \not\exists\beta\: (u,\beta)\in\sem{\hat I}{\hat q_2}
				}
				\uplus
				\mset{\left(u,\alpha\ominus\bigoplus_{\beta\mid (u,\beta)\in\sem{\hat
							I}{\hat q_2}}\beta\right)\mid (u,\alpha)\in\sem{\hat I}{\hat q_1},
				\exists\beta (u,\beta)\in\sem{\hat I}{\hat q_2}}                  \\
				       & =                &            &
				\mset{(u,\alpha)\mid (u,\alpha)\in \ansem{\hat I}{ q_1}, \not\exists\beta\: (u,\beta)\in\ansem{\hat I}{ q_2}
				}
				\uplus
				\mset{\left(u,\alpha\ominus\bigoplus_{\beta\mid
						(u,\beta)\in\ansem{\hat I}{ q_2}}\beta\right)\:\middle|\: (u,\alpha)\in\ansem{\hat I}{ q_1},
				\exists\beta (u,\beta)\in\ansem{\hat I}{ q_2}}                    \\
				       & =                &            &
				\mset{\left(u,\alpha\ominus\bigoplus_{\beta\mid
						(u,\beta)\in\ansem{\hat I}{ q_2}}\beta\right)\:\middle|\: (u,\alpha)\in\ansem{\hat I}{ q_1}
          } =
				     \ansem{\hat I}{q_1-q_2}.
			\end{align*}

% lax end
		\item \textbf{aggregation} for $k\in\bN$,
			$q\in\RA_k$, distinct $(i_j)_{1\leq j\leq k}$,
			terms~$t_1,\dots t_n$
			of max-index $\leq k$, monoid aggregate functions
			$f_1,\dots,f_n$,\\
      \hspace*{-5em}\begin{minipage}{\linewidth}
			\begin{align*}
				 & \phantom{{}={}}\sem{\hat
					I}{\widehat{\gamma_{i_1,\dots,i_m}[t_1:f_1,\dots,t_n:f_n](q)}}
				\\& =\sem{\hat I}{\gamma_{i_1,\dots,i_m}[t_1*\#(k+1):\hat
						f_1,\dots,t_n*\#(k+1):\hat
						f_n,\#(k+1):\delta\left(\bigoplus\right)](\hat q)}
				\\
				 & =\set{\left(v_1,\dots,v_m, \hat f_1\left(\mset{\scriptstyle t_1(u)*\alpha\mid (u,\alpha)\in\sem{\hat
					I}{\hat q},
				(u_{i_1},\dots,u_{i_m})=(v_1,\dots,v_m)}\right),\dots,\delta\left(
				\bigoplus_{\substack{(u,\beta)\in\sem{\hat I}{\hat q}                                      \\(u_{i_1},\dots,u_{i_m})=(v_1,\dots,v_m)}}\beta
				\right)\right)
				\:\middle|\:
				(v_1,\dots,v_m)\in\sem{\hat{I}}{\epsilon(\Pi_{\#i_1,\dots,\#i_m}(\hat q))}}                \\
				 & =\set{\left(v_1,\dots,v_m, \hat
				f_1\left(\mset{\scriptstyle t_1(u)*\alpha\mid (u,\alpha)\in\ansem{\hat
					I}{q},
				(u_{i_1},\dots,u_{i_m})=(v_1,\dots,v_m)}\right),\dots,\delta\left(
				\bigoplus_{\substack{(u,\beta)\in\sem{\hat I}{\hat q}                                      \\(u_{i_1},\dots,u_{i_m})=(v_1,\dots,v_m)}}\beta
				\right)\right)
				\:\middle|\:
				(u,\beta)\in \sem{\hat I}{\hat q}, (u_{i_1},\dots,
				u_{i_m})=(v_1,\dots,v_m)}                                                                  \\
				 & =\set{\left(v_1,\dots,v_m, \hat
				f_1\left(\mset{\scriptstyle t_1(u)*\alpha\mid (u,\alpha)\in\ansem{\hat
					I}{q},
				(u_{i_1},\dots,u_{i_m})=(v_1,\dots,v_m)}\right),\dots,\delta\left(
				\bigoplus_{\substack{(u,\beta)\in\ansem{\hat I}{q}                                      \\(u_{i_1},\dots,u_{i_m})=(v_1,\dots,v_m)}}\beta
				\right)\right)
				\:\middle|\:
				(u,\beta)\in \ansem{\hat I}{q}, (u_{i_1},\dots,
				u_{i_m})=(v_1,\dots,v_m)}                                                                  \\
				 & =\set{\left(v_1,\dots,v_m, \hat
				f_1\left(\mset{\scriptstyle t_1(u)*\alpha\mid (u,\alpha)\in\ansem{\hat
					I}{q},
				(u_{i_1},\dots,u_{i_m})=(v_1,\dots,v_m)}\right),\dots,\delta
				\left(
        \bigoplus_{(u,\beta)\in\ansem{\hat I}{\Pi_{\#i_1,\dots,\#i_m}(q)}}\beta
				\right)
				\right)
				\:\middle|\:
				(u,\beta)\in \ansem{\hat
				I}{\Pi_{\#i_1,\dots,\#i_m}(q)}}                                                            \\
				 & =\set{\left(v_1,\dots,v_m, \hat
				f_1\left(\mset{\scriptstyle t_1(u)*\alpha\mid (u,\alpha)\in\ansem{\hat
					I}{q},
				(u_{i_1},\dots,u_{i_m})=(v_1,\dots,v_m)}\right),\dots,\delta(\beta)\right)
				\:\middle|\:
				(u,\beta)\in \ansem{\hat
				I}{\epsilon(\Pi_{\#i_1,\dots,\#i_m}(q)}}                                                   \\
				 & =\ansem{\hat
         I}{\gamma_{i_1,\dots,i_m}[t_1:f_1,\dots,t_n:f_n](q)}. \qedhere
			\end{align*}
    \end{minipage}
	\end{itemize}
\end{proof}


\begin{example}\label{ex:rewriting}
	Returning to query $\qex$ from Example \ref{ex:semantics}, let us trace
	the rewriting rules applied by ProvSQL bottom up (for space reasons,
	$\mathit{Personnel}$ is abbreviated to $P$):
	\begin{align*}
		\qex={}                              & \epsilon(\Pi_{\#4}(P\bowtie_{\#4=\#8\land\#1<\#5}P))                                                              \\
		={}                                  & \epsilon(\Pi_{\#4}(\sigma_{\#4=\#8\land\#1<\#5}(P\times P)))                                                      \\
		\xrightarrow{(R\ref{rcp})}{}         & \epsilon(\Pi_{\#4}(\sigma_{\#4=\#8\land\#1<\#5}(\Pi_{\#1,\dots,\#4,\#6,\dots\#9,\#5\otimes\#10}(P\times P))))     \\
		\xrightarrow{(R\ref{rprojection})}{} & \epsilon(\Pi_{\#4,\#9}(\sigma_{\#4=\#8\land\#1<\#5}(\Pi_{\#1,\dots,\#4,\#6,\dots\#9,\#5\otimes\#10}(P\times P)))) \\
		\xrightarrow{(R\ref{rdupelim})}{}    &
		\gamma_{1}[\#2:\bigoplus](                                                                                                                               \\
		                                     & \quad\Pi_{\#4,\#9}(\sigma_{\#4=\#8\land\#1<\#5}(\Pi_{\#1,\dots,\#4,\#6,\dots\#9,\#5\otimes\#10}(P\times
		P))))
	\end{align*}
	If $\qexhat$ is this rewritten query, one can check that
	$\sem{\hat I}{\qexhat}=\ansem{\hat I}{\qex}$.
\end{example}
